\section{Introduction}

\modname{SelBiasEngine} (\pyfile{src/clenspy/selection/bsel.py}) implements
Costanzi et al.\ (2026)'s selection-affected halo bias: a cluster selected
at observed richness $\lob$ preferentially sits behind extra correlated
structure, so its two-halo term carries $\bsel(\theta)$ rather than the
ordinary $b(M,z)$. $\bsel(\theta)$ interpolates between two plateaus,
$\bsmall$ (the aperture, $\theta\lesssim\theta_\lambda$) and $\blarge$
(the field far outside it), from a closure of the projection operator
$\mathcal P[X]$. That closure and the fix derived this session are
documented in \pyfile{docs/selection\_bias.md} and
\pyfile{docs/plan-bsel-stable-closure.md}; this report is the validation
record, not a re-derivation.

This report uses two independent references, deliberately, because
agreement with only one would leave open whether a residual is a model
bug or a reference artefact:

\begin{itemize}
\item \textbf{Costanzi et al.\ (2026) Fig.~6} (\S\ref{sec:fig6}) --- the
  published model curve, digitized by hand. It carries no sampling
  noise, so any residual is attributable to a model or input mismatch,
  not measurement scatter.
\item \textbf{The Buzzard mock catalogue} (\S\ref{sec:mock}) --- the
  synthetic light-cone the published figure itself was built from, with
  real shot noise but a known generative cosmology and MOR (confirmed,
  not assumed, via \pyfile{costanzi\_notebook/} in the mock's own repo).
\end{itemize}

\S\ref{sec:methods} summarises the numerical methods shared by
$\bsel$ and $\Sprj$ without re-deriving them. \S\ref{sec:alpha} and
\S\ref{sec:cosmo} report the closure's sensitivity to the HOD slope
$\alpha$ and to $(\Omega_m,\sigma_8)$, respectively --- the systematic
uncertainty budget this validation implies for the model, not further
model-vs-reference checks.

All numbers in this report are generated by
\pyfile{papers/validations/scripts/*.py} from the same production code
under \pyfile{src/clenspy/}; none are hand-typed (see the Makefile).
